% 拉普拉斯算子（综述）
% license CCBYSA3
% type Wiki

本文根据 CC-BY-SA 协议转载翻译自维基百科\href{https://en.wikipedia.org/wiki/Laplace_operator}{相关文章}。

\subsection{拉普拉斯算子}
在数学中，拉普拉斯算子（Laplace operator 或 Laplacian）是一种微分算子，由欧几里得空间中标量函数的梯度的散度给出。它通常记作符号 $\nabla \cdot \nabla$、$\nabla^{2}$（其中 $\nabla$ 为纳布拉算子，nabla operator），或 $\Delta$。在笛卡尔坐标系中，拉普拉斯算子等于函数关于每个独立变量的二阶偏导数之和；在其他坐标系（如柱坐标系和球坐标系）中，拉普拉斯算子也有相应的有用形式。非正式地说，函数 $f$ 在点 $p$ 的拉普拉斯值 $\Delta f(p)$ 衡量的是：以 $p$ 为中心的小球上 $f$ 的平均值与 $f(p)$ 本身的偏离程度。

拉普拉斯算子以法国数学家皮埃尔–西蒙·拉普拉斯（Pierre-Simon de Laplace，1749–1827）命名。他最早将该算子应用于天体力学研究：给定质量密度分布下的引力势的拉普拉斯值是该密度分布的常数倍。满足拉普拉斯方程 $\Delta f = 0$ 的函数称为调和函数（harmonic functions），它们表示真空区域内可能的引力势。

拉普拉斯算子出现在许多描述物理现象的微分方程中：泊松方程（Poisson’s equation）描述电势与引力势，扩散方程（diffusion equation）描述热流与流体流动，波动方程（wave equation）描述波的传播，而薛定谔方程（Schrödinger equation）则描述量子力学中的波函数。在图像处理与计算机视觉领域中，拉普拉斯算子也被用于多种任务，如斑点检测（blob detection）与边缘检测（edge detection）拉普拉斯算子是最简单的椭圆型算子（elliptic operator），并且在霍奇理论（Hodge theory）及德拉姆上同调（de Rham cohomology）的结果中起着核心作用。
\subsection{定义}
拉普拉斯算子（Laplace operator）是定义在 $n$ 维欧几里得空间中的二阶微分算子，
其定义为梯度（$\nabla f$）的散度（$\nabla \cdot$）。
因此，若 $f$ 是一个二次可微的实值函数，则 $f$ 的拉普拉斯算子为：

\[
\Delta f = \nabla^{2} f = \nabla \cdot \nabla f,
\tag{1}
\]

其中符号 $\nabla$ 正式写作
\[
\nabla = \left( \frac{\partial}{\partial x_{1}}, \ldots, \frac{\partial}{\partial x_{n}} \right)。
\]

显式地说，在笛卡尔坐标系 $(x_{1}, \ldots, x_{n})$ 中，
函数 $f$ 的拉普拉斯算子是其关于各坐标的非混合二阶偏导数之和：

\[
\Delta f = \sum_{i=1}^{n} \frac{\partial^{2} f}{\partial x_{i}^{2}}。
\tag{2}
\]

作为一个二阶微分算子，拉普拉斯算子将 $C^{k}$ 类函数映射为 $C^{k-2}$ 类函数（当 $k \ge 2$ 时）。
它是一个线性算子 $\Delta : C^{k}(\mathbb{R}^{n}) \to C^{k-2}(\mathbb{R}^{n})$，
或更一般地，$\Delta : C^{k}(\Omega) \to C^{k-2}(\Omega)$，
其中 $\Omega \subseteq \mathbb{R}^{n}$ 为任意开集。

另一种等价定义为：
\[
\nabla^{2} f(\vec{x})
= \lim_{R \to 0} \frac{2n}{R^{2}} \left(f_{\text{shell}_{R}} - f(\vec{x})\right)
= \lim_{R \to 0} \frac{2n}{A_{n-1} R^{1+n}}
\int_{\text{shell}_{R}} \left( f(\vec{r}) - f(\vec{x}) \right) \, d r^{n-1},
\]
其中：
\begin{itemize}
  \item $n$ 是空间的维数；
  \item $f_{\text{shell}_{R}}$ 表示 $f$ 在半径为 $R$ 的 $n$ 维球面上的平均值；
  \item $\displaystyle \int_{\text{shell}_{R}} f(\vec{r}) \, d r^{n-1}$ 是 $f$ 在半径为 $R$ 的 $n$ 维球面上的面积积分；
  \item $A_{n-1}$ 是单位 $n$ 维球面的边界超体积（即单位球面的 $(n-1)$ 维面积）\cite{1}。
\end{itemize}
